\subsection{Непрерывные функции в точке (по Коши и по Гейне). Односторонняя непрерывность. Арифмитические свойства непрерывных функций.}

\subsubsection*{Определения непрерывности}

Пусть $f(x)$ определена в окрестности $U_a$.

\begin{tcolorbox}[title=Определение (через пределы)]
	Функция $f(x)$ называется \textbf{непрерывной в точке} $a$, если существует предел функции в этой точке и он равен значению функции:
	\[ \exists \lim_{x\to a} f(x) = f(a) \]
	
	\textit{Эквивалентная запись:} $\lim_{x\to a} f(x) = f(\lim_{x\to a} x)$. (Знак предела и функции можно менять местами).
\end{tcolorbox}

\begin{tcolorbox}[title=Определение по Коши ($\varepsilon - \delta$)]
	Функция $f(x)$ непрерывна в точке $a$, если:
	\[ \forall \varepsilon > 0\; \exists \delta_\varepsilon > 0:\; \forall x \in U_a \quad (|x - a| < \delta_\varepsilon \Rightarrow |f(x) - f(a)| < \varepsilon) \]
	\textit{Замечание:} В отличие от определения предела, здесь неравенство $|x-a| < \delta$ нестрогое ($0 <$ не нужно), так как при $x=a$ неравенство $|f(a)-f(a)|=0 < \varepsilon$ выполняется автоматически.
\end{tcolorbox}

\begin{tcolorbox}[title=Определение по Гейне (на языке последовательностей)]
	Функция $f(x)$ непрерывна в точке $a$, если для любой последовательности $\{x_n\} \subset U_a$, сходящейся к $a$ ($x_n \to a$), соответствующая последовательность значений функции сходится к $f(a)$:
	\[ \forall \{x_n\}: \lim_{n\to\infty} x_n = a \implies \lim_{n\to\infty} f(x_n) = f(a) \]
\end{tcolorbox}

\subsubsection*{Односторонняя непрерывность}

\begin{tcolorbox}
	\begin{itemize}
		\item $f(x)$ непрерывна в точке $a$ \textbf{справа}, если $\lim_{x\to a + 0} f(x) = f(a)$.
		\item $f(x)$ непрерывна в точке $a$ \textbf{слева}, если $\lim_{x\to a - 0} f(x) = f(a)$.
	\end{itemize}
\end{tcolorbox}
\textbf{Пример:} $y = [x]$ (целая часть).
В точке $x=1$: $f(1) = 1$.
$\lim_{x\to 1+0} [x] = 1$ (равен значению $\Rightarrow$ непрерывна справа).
$\lim_{x\to 1-0} [x] = 0$ (не равен значению).

\subsubsection*{Свойства непрерывных функций}

\textbf{Теорема (Арифметические свойства).}
Если $f(x)$ и $g(x)$ непрерывны в точке $a$, то функции:
\begin{enumerate}
	\item $f(x) \pm g(x)$
	\item $f(x) \cdot g(x)$
	\item $\frac{f(x)}{g(x)}$ (при условии $g(a) \neq 0$)
\end{enumerate}
также непрерывны в точке $a$.

\begin{tcolorbox}[title=Доказательство]
	Так как функции непрерывны, то $\lim_{x\to a} f(x) = f(a)$ и $\lim_{x\to a} g(x) = g(a)$.
	Используем свойства пределов:
	\begin{enumerate}
		\item $\lim (f(x) \pm g(x)) = \lim f(x) \pm \lim g(x) = f(a) \pm g(a)$.
		\item $\lim (f(x) \cdot g(x)) = \lim f(x) \cdot \lim g(x) = f(a) \cdot g(a)$.
		\item $\lim \frac{f(x)}{g(x)} = \frac{\lim f(x)}{\lim g(x)} = \frac{f(a)}{g(a)}$.
	\end{enumerate}
	Значения пределов совпадают со значениями выражений в точке $a$, что и означает непрерывность.
\end{tcolorbox}

